\documentclass[11pt]{article}
\usepackage[utf8]{inputenc}
\usepackage{amsmath,amssymb}
\usepackage[margin=1in]{geometry}
\usepackage{hyperref}
\emergencystretch=3em

\title{Taylor-data identification of the double-series coefficients}
\author{Claude, with Daniel Murfet}
\date{2026-09-07}

\begin{document}
\maketitle
\sloppy

\begin{abstract}
Astra~\#10's R7. The Taylor-tree theorems take coefficient arrays $x,y\colon\mathbb N^2\to\mathbb R$
as input; the paper writes them as Taylor coefficients of the phase and amplitude. We prove the
identification: if $\xi(u,v)=\sum_{ij}x_{ij}u^iv^j$ on the open ball $|u|,|v|<\rho$ with
$\|x\|_\rho<\infty$, then $\partial_v^j\partial_u^i\xi(0,0)=i!\,j!\,x_{ij}$
(\texttt{taylorData\_of\_eq\_dblSum}), so the coefficients are determined by the function on any
ball (\texttt{dblSum\_injective}). The one-variable engine is \texttt{iteratedDeriv\_zero\_of\_hasSum}:
$\partial^nf(0)=n!\,c_n$ for $f(y)=\sum c_ny^n$ on $|y|<r$ with $\sum|c_n|r^n<\infty$, a direct
consequence of Mathlib's \texttt{HasFPowerSeriesOnBall.factorial\_smul} applied to the scalar series
\texttt{FormalMultilinearSeries.ofScalars}. Zero \texttt{sorry}/\texttt{axiom}.
\end{abstract}

\section{The things formalised}
{\small
\begin{verbatim}
theorem iteratedDeriv_zero_of_hasSum (c : ℕ → ℝ) (r) (hr : 0 < r)
    (hs : Summable fun n => |c n| * r ^ n) (f : ℝ → ℝ)
    (hf : ∀ y, |y| < r → HasSum (fun n => c n * y ^ n) (f y)) (n) :
    iteratedDeriv n f 0 = n.factorial * c n
def rowNorm (ρ) (x) (i) : ℝ := ∑' j, |x (i, j)| * ρ ^ j
theorem row_summable / rowNorm_mul_summable / abs_faceU_le_rowNorm / faceU_hasSum
theorem dblSum_hasSum_faceU (ρ) (hρ) (x) (hx : WSummable ρ x) (u v)
    (hu : |u| ≤ ρ) (hv : |v| ≤ ρ) :
    HasSum (fun i => faceU x i v * u ^ i) (dblSum x u v)
theorem taylorData_of_eq_dblSum (ρ) (hρ) (x) (hx) (ξ : ℝ → ℝ → ℝ)
    (hξ : ∀ u v, |u| < ρ → |v| < ρ → ξ u v = dblSum x u v) (i j) :
    iteratedDeriv j (fun v => iteratedDeriv i (fun u => ξ u v) 0) 0
      = i.factorial * j.factorial * x (i, j)
theorem dblSum_taylorData ... (ξ = dblSum x)
theorem dblSum_injective (ρ) (hρ) (x x') (hx hx')
    (h : ∀ u v, |u| < ρ → |v| < ρ → dblSum x u v = dblSum x' u v) : x = x'
\end{verbatim}
}

\section{Mathematics}
For a scalar power series $\sum c_ny^n$ with $\sum|c_n|r^n<\infty$, the formal multilinear series
$p_n=c_n\cdot\mathrm{mkPiAlgebraFin}$ has radius $\ge r$ and sums to $f$ on the ball, so
\texttt{HasFPowerSeriesOnBall f p 0 r}; Mathlib's \texttt{factorial\_smul} then gives
$n!\,p_n(1,\dots,1)=D^nf(0)(1,\dots,1)=\partial^nf(0)$, and $p_n(1,\dots,1)=c_n$. Nothing about
the function off the ball is used, which is what makes the lemma applicable to functions that only
agree with the series near the origin. The two-variable statement is two applications: for fixed
$|v|<\rho$ the function $u\mapsto\xi(u,v)$ is the series $\sum_ia_i(v)u^i$ with
$a_i(v)=\sum_jx_{ij}v^j$ (\texttt{faceU}), absolutely convergent since
$|a_i(v)|\le\sum_j|x_{ij}|\rho^j$ and $\sum_i(\sum_j|x_{ij}|\rho^j)\rho^i=\|x\|_\rho$; so
$\partial_u^i\xi(0,v)=i!\,a_i(v)$. Then $v\mapsto i!\,a_i(v)$ is the series with coefficients
$i!\,x_{ij}$, and one more application gives $i!\,j!\,x_{ij}$. Injectivity follows by applying the
identification to $x'$ against the function $\mathrm{dblSum}\,x$.

\section{Paper-facing meaning}
This closes the gap between the paper's phrasing (``$x_{ij}$ are the Taylor coefficients of $\xi$
at the origin'') and the Lean input (an abstract array with $\|x\|_\rho<\infty$): any analytic
$\xi$ on the $\rho$-ball with such a representation has exactly these coefficients, and the
canonical Taylor-tree coefficients $A_\alpha,B_\alpha$ are therefore functions of the germ of
$(\xi,\eta)$ at the origin. With units~138--143 the whole chain germ $\mapsto$ coefficient arrays
$\mapsto$ amplitude coefficients $\mapsto(A_\alpha,B_\alpha)$ is now formalised, the last step
locally Lipschitz. Astra~\#10's R7 is complete; what remains from that consult is the R9 expectation
bridge for $\mathbb E[\log Z]$, which needs the §4.3 fluctuation-function definition first.

\section{Lean notes}
\begin{itemize}
\item \texttt{FormalMultilinearSeries.le\_radius\_of\_summable\_norm} and
\texttt{ofScalars\_norm} (needs \texttt{NormOneClass}, fine over $\mathbb R$) give the radius
bound; \texttt{ofScalars\_apply\_eq} evaluates $p_n(y,\dots,y)=c_n\cdot y^n$;
\texttt{iteratedDeriv\_eq\_iteratedFDeriv} is \texttt{rfl} and bridges to \texttt{factorial\_smul}.
\item Ball membership for an \texttt{NNReal} radius: \texttt{Metric.eball\_coe}
(\texttt{Metric.emetric\_ball\_nnreal} is deprecated on this pin), then
\texttt{Metric.mem\_ball, dist\_zero\_right}.
\item \texttt{summable\_prod\_of\_nonneg} with a \texttt{by positivity} nonnegativity argument fails
when the term is presented as an unreduced lambda application; pass an explicit
\texttt{mul\_nonneg (abs\_nonneg \_) (pow\_nonneg \_ \_)} term instead.
\item The style linter flags goal-changing \texttt{show} (use \texttt{change}); the project-wide
``Copyright too short!'' header lint fires on every Grammar file and is not specific to this unit.
\item Row-wise resummation is \texttt{HasSum.prod\_fiberwise} applied to the absolutely convergent
double family, with each row's \texttt{HasSum} from \texttt{faceU\_hasSum.mul\_right}.
\end{itemize}

\end{document}
